\begin{diagram}[
  type=flow,
  label={fig:sec02-ingestion-semantics},
  caption={Every event on this path is checkpointed, retried, and deduplicated before it reaches the sink; that sequence is what makes the sink safe to call idempotent under at-least-once delivery.},
  source={Conceptual diagram.},
  description={A top-down flow of eight stages in reading order: Source, Extractor or Producer, Broker or Buffer, Raw Landing, Checkpoint, Retry, Deduplication, and Idempotent Sink. Every event passes through all eight stages in sequence; the checkpoint, retry, and deduplication stages exist because at-least-once delivery can create duplicate or replayed events, so the flow deduplicates before writing to a sink built to tolerate repeated writes.},
  minspace={0.62\textheight}
]
% Eight stages are too dense for the default horizontal spacing, so this
% diagram places them directly with \RKNode/\flowedge in a vertical column.
\RKNode{source}{0}{0}{Source}
\RKNode{extractor}{0}{-1.85}{Extractor or Producer}
\RKNode{broker}{0}{-3.70}{Broker or Buffer}
\RKNode{landing}{0}{-5.55}{Raw Landing}
\RKNode{checkpoint}{0}{-7.40}{Checkpoint}
\RKNode{retry}{0}{-9.25}{Retry}
\RKNode{dedup}{0}{-11.10}{Deduplication}
\RKNode{sink}{0}{-12.95}{Idempotent Sink}
\flowedge{source}{extractor}
\flowedge{extractor}{broker}
\flowedge{broker}{landing}
\flowedge{landing}{checkpoint}
\flowedge{checkpoint}{retry}
\flowedge{retry}{dedup}
\flowedge{dedup}{sink}
\end{diagram}
